\documentclass[10pt]{article}

\usepackage[margin=0.72in,headheight=21pt]{geometry}
\usepackage{amsmath}
\usepackage{array}
\usepackage{booktabs}
\usepackage{enumitem}
\usepackage{fancyhdr}
\usepackage{graphicx}
\usepackage{hyperref}
\usepackage{microtype}
\usepackage{tabularx}
\usepackage[table]{xcolor}

\definecolor{AuroraBlue}{HTML}{246BCE}
\definecolor{AuroraLight}{HTML}{EAF3FF}
\definecolor{RuleGray}{HTML}{C9CED6}
\definecolor{TextGray}{HTML}{4E5661}
\definecolor{GoodGreen}{HTML}{18794E}
\definecolor{WarnRed}{HTML}{B42318}

\hypersetup{
  colorlinks=true,
  linkcolor=AuroraBlue,
  urlcolor=AuroraBlue,
  pdfauthor={AURORA Risk Engine Team},
  pdftitle={Risk Engine Model Evaluation}
}
\urlstyle{same}

\pagestyle{fancy}
\fancyhf{}
\lhead{\textcolor{TextGray}{AURORA Risk Engine}}
\rhead{\textcolor{TextGray}{Five-session volatility}}
\cfoot{\thepage}
\renewcommand{\headrulewidth}{0.4pt}
\renewcommand{\headrule}{\hbox to\headwidth{\color{RuleGray}\leaders\hrule height \headrulewidth\hfill}}

\setlength{\parindent}{0pt}
\setlength{\parskip}{5pt}
\setlist[itemize]{leftmargin=1.3em,itemsep=2pt,topsep=3pt}
\setlist[enumerate]{leftmargin=1.5em,itemsep=2pt,topsep=3pt}
\renewcommand{\arraystretch}{1.18}

\newcommand{\metric}[1]{\textbf{\textcolor{AuroraBlue}{#1}}}
\newcommand{\positive}[1]{\textbf{\textcolor{GoodGreen}{#1}}}
\newcommand{\negative}[1]{\textbf{\textcolor{WarnRed}{#1}}}

\begin{document}

{\LARGE\bfseries Risk Engine Model Evaluation}

\vspace{0.25em}
{\large Five-session realized-volatility forecasting}

\vspace{0.75em}
\colorbox{AuroraLight}{%
  \parbox{\dimexpr\textwidth-2\fboxsep\relax}{%
    \textbf{Decision.}
    Keep \textbf{HAR-X + News} as the formal Risk Engine.
    Price-only XGBoost is the strongest point-forecast challenger, while
    XGBoost with a constrained news-attention residual is the best
    news-aware research challenger. Neither XGBoost variant is ready to
    replace the formal model.
  }%
}

\section{Key findings}

\begin{itemize}
  \item \textbf{XGBoost Gamma has the lowest out-of-fold QLIKE:}
        \metric{0.6571}, a \metric{10.31\%} improvement over HAR-X
        (\(p=0.0010\)).
  \item \textbf{HAR-X + News is the best production compromise:}
        QLIKE improves by \metric{5.24\%}, the news increment is significant
        (\(p=0.0011\)), and the realized-to-forecast calibration ratio is
        \metric{0.981}.
  \item \textbf{Residual MLP does not establish a reliable advantage:}
        average QLIKE improves by 4.64\%, but the date-aggregated DM test is
        not significant (\(p=0.915\)).
  \item \textbf{News can improve XGBoost only under a constrained design:}
        an attention-based Gamma residual reduces XGBoost QLIKE by
        \positive{1.30\%}, but worsens high-volatility QLIKE by
        \negative{4.17\%}.
  \item \textbf{FinBERT sentiment does not improve the XGBoost residual model}
        in the current panel. Its measurable XGBoost news contribution comes
        from attention, recency, and causal coverage states.
\end{itemize}

\section{Evaluation protocol}

\begin{tabularx}{\textwidth}{>{\bfseries}p{0.26\textwidth}X}
\toprule
Target & Future five-session daily realized volatility; no price-direction label is used. \\
Universe & 21 FNSPID research symbols; 27,027 stock-date outer-fold forecasts. \\
Training history & 9 April 2013 to 28 December 2023. \\
Outer evaluation & Annual expanding walk-forward tests for 2018--2023. \\
Leakage control & Five-session embargo at each outer boundary; model, feature, and calibration decisions use prior data only. \\
Weighting & Stock-equal loss weighting so high-news-volume symbols cannot dominate. \\
Primary metric & QLIKE on realized variance; lower is better. \\
Supporting evidence & Log-RMSE, log-\(R^2\), calibration ratio, high-volatility QLIKE, per-stock gains, Newey-West DM tests, and moving-block bootstrap intervals. \\
\bottomrule
\end{tabularx}

\section{Models}

\subsection{HAR-X}

HAR-X is an interpretable pooled log-volatility model. The formal price
specification uses six causal features:
\[
\log \widehat{\sigma}^{\,P}_{t,5}
=
\beta_0
+\beta_1\log RV_{t,5}
+\beta_2\log RV_{t,22}
+\beta_3\log RV_{t,66}
+\beta_4\log PARK_{t,5}
+\beta_5\log PARK_{t,22}
+\beta_6\log |r_t|.
\]
The fitted log forecast is exponentiated and corrected with a smearing factor.
HAR-X provides a smooth and interpretable reference, but misses more of the
high-volatility regime than the nonlinear challengers.

\subsection{HAR-X + News}

The formal five-session model combines the HAR-X price forecast with a
regularized Gamma variance-ratio term:
\[
\widehat{\sigma}_{t,5}
=
\widehat{\sigma}^{\,P}_{t,5}
\times
\exp\!\left[
  \frac{1}{2}\gamma
  \frac{\log(1+N_t)}{s_N}
\right].
\]
Here \(N_t\) is the causal news count and \(s_N\) is the training-period scale.
A zero-news observation is retained with \(N_t=0\); it is not dropped or
replaced by a sentiment-neutral article. The current deployable news term uses
\texttt{log\_count}, selected in all six outer folds.

\subsection{XGBoost Gamma}

The XGBoost challenger uses a Gamma objective to predict positive realized
variance from 34 price features. The inputs include HAR features, OHLC
range estimators, semivariances, jump and extreme-return measures, volume,
liquidity, SPY market risk, and cross-sectional risk states. For each outer
year, shallow-tree hyperparameters are chosen from the original 96-configuration
inner search, with early stopping and an internally selected HAR blend.

\subsection{Residual MLP}

The neural candidate uses 68 inputs: 34 price features and 34 deployable
news/state features. Three regularized architectures are considered inside
each outer training window:
\[
64\!-\!32,\qquad 128\!-\!64,\qquad 128\!-\!64\!-\!32.
\]
The networks use LayerNorm, GELU activations, dropout, AdamW, and a bounded
log-volatility residual around a fold-local HAR forecast. The candidate has
6,721--19,649 trainable parameters depending on the selected architecture.

\section{Core out-of-fold results}

\begin{table}[ht]
\centering
\caption{Five-session point-forecast performance, 2018--2023 OOF.
Lower QLIKE and log-RMSE are better.}
\small
\begin{tabular}{lrrrrrrr}
\toprule
Model & QLIKE & Gain & Log-RMSE & Log-\(R^2\) & Calib. & High-vol & DM \(p\) \\
\midrule
XGBoost Gamma
  & \textbf{0.6571} & \textbf{10.31\%} & 0.5685 & 0.291 & 0.892 & \textbf{1.3539} & 0.0010 \\
HAR-X + News
  & 0.6943 & 5.24\% & 0.5400 & 0.360 & \textbf{0.981} & 1.6808 & 0.0011 \\
Residual MLP
  & 0.6987 & 4.64\% & 0.5567 & 0.320 & 0.917 & 1.5654 & 0.9151 \\
HAR-X
  & 0.7327 & Reference & \textbf{0.5192} & \textbf{0.408} & 1.057 & 1.9986 & -- \\
\bottomrule
\end{tabular}
\end{table}

\textbf{Interpretation.}
XGBoost is the strongest model under the primary QLIKE objective and produces
the lowest high-volatility loss. HAR-X has better ordinary log-error metrics,
which reflects smoother performance in normal regimes, but it underperforms on
variance-ratio loss. HAR-X + News gives up some point-forecast performance
relative to XGBoost in exchange for a near-unit calibration ratio, validated
news contribution, simple serialization, and stable zero-news behavior.

The MLP improves in five of six annual folds and 81.0\% of stocks, but the
daily loss-difference series is unstable enough that its aggregate DM
\(p\)-value remains 0.915. It is not a promotion candidate.

\section{Validated news contribution}

\subsection{Formal HAR-X + News}

\begin{table}[ht]
\centering
\caption{News increment and tail-risk validation for the formal model.}
\begin{tabular}{lr}
\toprule
Metric & Result \\
\midrule
Stock-equal QLIKE gain vs HAR-X & \textbf{5.24\%} \\
95\% moving-block bootstrap interval & \textbf{+2.38\% to +8.83\%} \\
DM test & \(p=0.001087\) \\
Positive outer years & 5/6 \\
Positive stocks & 85.7\% \\
Median stock gain & 7.51\% \\
High-volatility QLIKE gain & 13.56\% \\
VaR-95 breach rate & 4.32\% \\
95\% risk-band coverage & 96.34\% \\
Expected-shortfall ratio & 0.983 \\
\bottomrule
\end{tabular}
\end{table}

The formal news model passes the historical accuracy, breadth, significance,
and tail-calibration gates. Its only unresolved evidence requirement is a
mature live-RSS evaluation period.

\subsection{Can news improve XGBoost?}

Two architectures were tested without removing zero-news samples.

\textbf{Direct concatenation failed.}
Appending news fields to the 34 price inputs increased QLIKE from 0.6571 to
0.6855, a \negative{4.32\% deterioration}. News-active rows deteriorated by
2.73\%, high-volatility rows by 7.93\%, and news features represented only
1.67\% of final-tree total gain.

\textbf{A constrained residual layer found a smaller positive signal.}
The price-only XGBoost forecast was held fixed and news was used only to model
its remaining variance ratio:
\[
\widehat{\sigma}_{t,5}
=
\widehat{\sigma}^{\,XGB}_{t,5}
\times
GammaRatio(\text{news}_{t}).
\]

\begin{table}[ht]
\centering
\caption{News residual families on top of price-only XGBoost.}
\scriptsize
\begin{tabular}{lrrrrrrr}
\toprule
News family & QLIKE & Gain & DM \(p\) & Bootstrap 95\% & Years & Stocks & High-vol \\
\midrule
Attention
  & \textbf{0.6486} & \textbf{+1.30\%} & 0.00137 & +0.51\% to +2.67\% & 4/6 & 81.0\% & \negative{-4.17\%} \\
Log count
  & 0.6531 & +0.61\% & 0.00027 & +0.27\% to +1.39\% & 5/6 & 71.4\% & -2.58\% \\
All deployable
  & 0.6518 & +0.81\% & 0.05235 & -0.10\% to +2.44\% & 4/6 & 71.4\% & -4.58\% \\
Attention + FinBERT
  & 0.6627 & -0.84\% & 0.64580 & -3.12\% to +1.89\% & 4/6 & 57.1\% & -8.12\% \\
\bottomrule
\end{tabular}
\end{table}

The attention family contains article frequency, short-window counts, time
since the latest article, and causal coverage/silence states. Its average gain
remains significant after a conservative four-family multiple-test correction,
but it fails the high-volatility non-degradation gate. FinBERT polarity and
distribution fields do not add value to this XGBoost residual specification.

\section{External generalization and production status}

The independent 2024--2026 price-backbone test contains 40 symbols and 25,400
mature five-session observations.

\begin{table}[ht]
\centering
\caption{HAR-X price-backbone external validation against trailing-22-day volatility.}
\small
\begin{tabular}{lrrrrr}
\toprule
Scope & Symbols & HAR QLIKE & Baseline QLIKE & Gain & DM \(p\) \\
\midrule
All external panel & 40 & 0.5829 & 0.6354 & 8.26\% & 0.00345 \\
Original research symbols & 21 & 0.6661 & 0.7195 & 7.42\% & 0.02742 \\
Unseen symbols & 19 & 0.4910 & 0.5425 & 9.48\% & 0.00626 \\
Current six-symbol portfolio & 6 & 0.6181 & 0.6262 & 1.30\% & 0.82950 \\
\bottomrule
\end{tabular}
\end{table}

Across all 40 symbols, the VaR-95 breach rate is 4.60\%, the 95\% risk-band
coverage is 95.37\%, and the expected-shortfall ratio is 0.914. The unseen
scope is mostly broad and sector ETFs, so it is evidence of cross-symbol
transfer rather than a complete unseen-single-stock test.

\newpage
The live news contract remains the final production gap. FNSPID validates the
historical news increment, but the system still needs at least 60 mature
five-session RSS forecasts to confirm that live source intensity and coverage
transfer correctly.

\section{Final model decision}

\begin{enumerate}
  \item \textbf{Formal output: HAR-X + News.}
        It is the only candidate that combines a statistically significant news
        contribution with acceptable VaR, risk-band, expected-shortfall, and
        high-volatility behavior.
  \item \textbf{Primary challenger: price-only XGBoost Gamma.}
        It has the best OOF QLIKE, but needs external unseen-stock validation,
        component-specific tail calibration, and a news specification that does
        not degrade high-volatility performance.
  \item \textbf{News-aware XGBoost research ceiling: attention residual.}
        The +1.30\% average gain is real, but the -4.17\% high-volatility result
        blocks promotion.
  \item \textbf{Residual MLP: research only.}
        More cross-sectional data and a longer live-news history are required
        before neural complexity can be justified.
\end{enumerate}

\section*{Reproducibility artifacts}

\small
\begin{itemize}
  \item Core comparison:
        \path{reports/risk_engine_presentation/model_comparison.csv}
  \item Annual OOF results:
        \path{reports/risk_engine_presentation/yearly_qlike.csv}
  \item XGBoost news-family comparison:
        \path{reports/risk_engine_xgb_news/overlay_family_comparison.csv}
  \item Formal promotion gates:
        \path{reports/risk_engine_optimization/promotion_gates.json}
  \item Current external validation:
        \path{data/processed/current_risk_test/current_validation_summary.csv}
  \item Formal checkpoint:
        \path{data/processed/risk_model.json}
\end{itemize}

\end{document}
